{\small\opmaakproza
\begin{enumerate}[leftmargin=1.5em,itemsep=3pt,label=\textcolor{kleur1}{\sffamily\bfseries\arabic*}]
\item \textbf{Lees de inleiding en de figuren goed.} Gegevens die je bij een latere vraag nodig hebt, staan vaak
  al bovenaan de opgave, in een tabel of in een grafiek.
\item \textbf{Schrijf de formule op.} Noteer eerst de formule die je gebruikt, vul dan de waarden in en reken pas daarna.
  Er is meestal een punt voor het gebruik van de juiste formule.
\item \textbf{Zoek op in je informatieboek.} Constanten, stofeigenschappen en formules staan in Binas of ScienceData.
  Noem waar je een waarde vandaan haalt als die niet in de opgave staat.
\item \textbf{Vergeet de eenheid niet.} Zonder eenheid, of met een eenheid die niet bij de grootheid past, krijg je het
  punt voor het completeren van de berekening niet. Reken in SI-eenheden: ms in s, km/h in m/s, kWh in J.
\item \textbf{Let op significante cijfers.} Geef de uitkomst in het aantal cijfers dat past bij de gegevens.
  Rond pas aan het eind af en reken tussendoor met minstens zoveel cijfers.
\item \textbf{Lees grafieken nauwkeurig af.} Teken bij een raaklijn of een oppervlakte onder de grafiek je hulplijnen
  in de figuur of op de uitwerkbijlage, zodat de corrector ziet wat je hebt gedaan.
\item \textbf{Bij \emph{leg uit} of \emph{beredeneer}} noem je de natuurkundige wet of formule, en zeg je wat er met
  de grootheden gebeurt. Een los antwoord als \emph{groter} of \emph{kleiner} levert meestal geen punten op.
\item \textbf{Bij \emph{toon aan}} werk je zichtbaar naar de gegeven waarde toe en geef je je uitkomst in meer cijfers
  dan de gegeven waarde.
\item \textbf{Kijk naar de punten.} Een vraag van 4 punten vraagt meer stappen dan een vraag van 2. Houd ongeveer
  2 minuten per punt aan en ga door als je vastloopt.
\item \textbf{Controleer je antwoord}: is de orde van grootte logisch? Een fietser van 300 m/s of een lamp van 5 MW klopt niet.
\end{enumerate}\par}
